% 【本文件写什么】api-v0 契约层，叶子，只写语义不写实现
% - crate 路径：os/components/wateros-vfs/vfs-api/api-v0/src/
% - 列出并说明：pub trait、核心类型、错误枚举、常量
% - 每个公开 API 的前置条件、后置条件、线程/中断上下文限制
% - 与 Linux/用户态约定的对齐点与 deliberate 差异
% - v0 版本当前行为与后续可替换 extension 点
% 【事实来源】
% - 上述 crate 下全部 .rs 源文件
% - docs/exports/public-api/ 中对应符号表，以源码为准
% 【不写什么】
% - impl 内部数据结构、汇编、具体算法步骤

\subsubsection{api-v0}

\paragraph{模块划分。}
\code{error}，即 \code{VfsError}/\code{VfsResult}；\code{kind}，即 \code{VfsFsKind}、\code{VfsCapability}；\code{meta}，即 \code{VfsMetadata}、\code{VfsDirEntry}；\code{path}，即 \code{normalize\_absolute\_path}、\code{validate\_root\_file\_name}；\code{resolve}，即 \code{resolve\_open\_path}、\code{resolve\_against\_cwd}；\code{handle}，即 \code{VfsOpenOps}、\code{VfsIoHandle}、\code{VfsOpenFlags}；\code{fd}，即 \code{VfsFdSession}；\code{mount}，即 \code{VfsMountOps}、\code{RootRwSession}；\code{namespace}，即 \code{VfsMountTable}；\code{dev}，即 \code{VfsDevInventory}；\code{root\_read}，即 \code{SingleRootReadView}；\code{backend}，即 \code{VfsBackend} 组合 trait。

\paragraph{\code{VfsBackend}。}
后端须同时实现\code{SingleRootReadView}、\code{VfsMountOps}、\code{VfsDevInventory}、\code{VfsOpenOps}、\code{VfsMountTable}。本crate不依赖\code{wateros-fs}；跨FS桥接仅在\code{impl-fs-bridge}。

\paragraph{路径契约。}
\code{normalize\_absolute\_path}折叠\code{//}、\code{.}、\code{..}，结果以\code{/}开头。\code{resolve\_open\_path}可经\code{register\_open\_path\_resolver}挂接per-task cwd解析，由 \code{impl-fd-session} 在启动期注册。\code{VfsOpenFlags}为语义层标志位，含 \code{READ}、\code{WRITE}、\code{CREATE}、\code{TRUNC}、\code{APPEND}、\code{DIRECTORY}，与Linux \code{O\_*}逐步对齐。

\paragraph{\code{VfsIoHandle}与\code{VfsFdSession}。}
\code{VfsIoHandle}描述已打开对象的\code{read}/\code{write}/\code{seek}/\code{read\_at}/\code{write\_at}/\code{metadata}/\code{close}；默认未覆盖方法返回\code{Unsupported}。\code{VfsFdSession}管理per-task fd表：\code{get\_io}、\code{alloc\_fd}、\code{close\_fd}；trait默认\code{alloc\_fd}为\code{Unsupported}，真实分配在\code{impl-fd-session}。

\paragraph{错误映射。}
\code{VfsError}与\code{FsError}一一对应，bridge 内经 \code{map\_fs\_err} 映射，syscall层再映射为\code{-errno}。刻意差异：\code{VfsMetadata}含\code{mount\_id}、\code{device\_major/minor}以支持\code{stat}与挂载表查询。

\paragraph{上下文限制。}
\code{no\_std}契约；\code{resolve}模块含\code{static mut}解析钩子，假定启动期单次注册。fd/cwd会话状态与\code{task::TaskId}生命周期绑定，由impl保证fork/clone继承语义。
